\subsection{Delta Learning: Current Understanding and Open Questions}
\label{sec:delta-learning}

The target transformations discussed above correspond mathematically to what is commonly referred to as delta learning. Rather than learning a reference property directly, delta learning defines a new prediction task by subtracting the prediction of a complementary baseline model,
\[
  \Delta y = y - b,
\]
where \(y\) is the reference target and \(b\) is a lower-fidelity baseline. The machine-learning model therefore learns only the residual, while the final prediction is reconstructed by combining the baseline and residual models. This idea was formalized in quantum chemistry by Ramakrishnan, Dral, Rupp, and von Lilienfeld, who showed that machine-learning corrections to inexpensive quantum-chemical methods could recover higher-level accuracy at reduced computational cost and with enhanced learning efficiency \cite{ramakrishnan2015delta}. Since then, delta learning has been applied to semiempirical-to-DFT corrections, DFT-to-coupled-cluster corrections, potential-energy surfaces, excited-state properties, solution-phase molecular properties, and machine-learning interatomic potentials, where it is especially useful in data-scarce regimes \cite{atz2022deltaqml,chen2023solutiondelta,dral2020hierarchical,krug2025aha}.

Despite the reported benefits, and sometimes reported failures, of delta learning, the mechanism by which it improves learning remains unclear. The literature suggests three broad explanations. The first is the \emph{variance-reduction hypothesis}: the residual target has lower variance than the original reference, making optimization easier. The second is the \emph{target-correlation hypothesis}: the baseline is strongly correlated with the reference target, so the residual is expected to be easier to learn. The third is the \emph{physical-prior hypothesis}: the baseline already contains physically meaningful information, leaving the model to learn only the missing correction \cite{ramakrishnan2015delta,atz2022deltaqml}. Each explanation is plausible and supported by empirical observations, but each remains largely descriptive rather than predictive. They explain successful applications after training, yet provide little guidance for selecting baselines before training.

This distinction is important because the ultimate goal is to establish delta-learning as a predictive target design principle. If variance reduction alone were sufficient, any baseline that reduced target variance should improve learning. If baseline accuracy alone determined success, increasingly accurate baselines should always produce increasingly transferable residuals. Existing studies, however, suggest that neither interpretation is universally valid \cite{atz2022deltaqml,chen2023solutiondelta,krug2025aha}. Similarly, describing a residual as ``smoother'' merely restates the observed improvement unless smoothness itself is defined with respect to the representation space of the learning model.

The central question therefore remains unresolved: for a fixed model architecture, what properties should a baseline possess to produce a target that is intrinsically more learnable? In this work, we investigate this question by constructing baselines with different physical origins, predictive accuracies, and statistical properties, and evaluating how each reshapes the learning problem. Rather than treating delta learning as an empirical recipe, we seek initial design principles that connect baseline choice, target landscape, and model transferability.
